%!TEX root = ../AdvMathLecture.tex
%
% Section: Path Independence
% 

%%~~~~~~~~~~~~~~~~~~~~~~   []   ~~~~~~~~~~~~~~~~~~~~~~~~~~~
%\begin{frame}
%\frametitle{\Large {\color[rgb]{0.02,0.20,1.0} \UL{Path Independence}}}
%\scriptsize
%
%\begin{block}{Key Concepts:}
%%-------LIST---------
%\begin{itemize}
%	\item[...] Path Independence
%	\item[...] Conservative Vector Field
%	\item[...] Potential Functions
%\end{itemize}
%%-------------------
%\end{block}
%
%\vspace{1.2cm}
%
%\begin{block}{Mathematical Tools:}
%%-------LIST---------
%\begin{itemize}
%	\item[...] Compute Line Integrals.
%	\item[...] Determine potential functions corresponding to conservative vector fields.
%	\item[...] Determine whether a field is conservative.
%\end{itemize}
%%-------------------
%\end{block}
%
%%\vspace{3.0cm}
%\end{frame}

%~~~~~~~~~~~~~~~~~~~~~~~~~~~~~~~~~~~~~~~~~~~~~~~~~~~~~~~~~
\begin{frame}
\frametitle{Topic: Path Independence} 
\begin{columns}\begin{column}{12.0cm}

{I. Introduction}

\vspace{0.1cm}

\scriptsize

Question: Does the amount of work required to move a mass from $A$ to $B$ in a force field depend on the path you take?

\begin{center}
	\includegraphics[scale = 0.18]{Figures/10_2_Fig1.png} \hspace{0.1cm} \includegraphics[scale = 0.18]{Figures/10_2_Fig2.png}
\end{center}

Answer: \ON<2->{\B{It depends on the force field.}}
\end{column}\begin{column}{0.0cm}\end{column}\end{columns}
\end{frame}

%%%~~~~~~~~~~~~~~~~~~~~~~~~~~~~~~~~~~~~~~~~~~~~~~~~~~~~~~~~~
%%\begin{frame}
%%\frametitle{}
%%\scriptsize
%%\begin{columns}
%%\begin{column}{12.0cm}
%%
%%\EX Let $C$ be a curve from point $(0,0,0)$ to $(1,1,0)$.  If \, ${\bf F} = [0,xy,0]$, show that the line integral $\WORK$ is path dependent.
%%
%%
%%\SOL $\{$WTS: there are curves, $C_1$ and $C_2$, such that \quad $\D{\int \limits_{C_1} \FDR \ne \int \limits_{C_2} \FDR }$ $\}$
%%
%%\hspace{8.8cm} \includegraphics[scale = 0.18]{Figures/10_2_Fig3.png}
%%
%%\vspace{6.2cm}
%%
%%\begin{columns}
%%\begin{column}{5.0cm}
%%	
%%\end{column}
%%\begin{column}{5.0cm}
%%	
%%\end{column}
%%\end{columns}
%%
%%\end{column}
%%\begin{column}{0.0cm}\end{column}
%%\end{columns}
%%\end{frame}

%~~~~~~~~~~~~~~~~~~~~~~~~~~~~~~~~~~~~~~~~~~~~~~~~~~~~~~~~~
\begin{frame}
\begin{columns}\begin{column}{12.0cm}

{II. Path Independence}
\scriptsize
\begin{center}
\setlength{\extrarowheight}{-0.2cm}
\begin{tabular}{|cp{10cm}|}
	\hline
	\hline
	&\\
	\DEF & Let $C_1$ be a curve from point $A$ to $B$. The line integral $\WORK$ is said to be {\bf path independent} if
	\ON<2->{\B{for any other curve, $C_2$, between $A$ and $B$}}\\
	&\\
	\multicolumn{2}{|c|}{$\D{
	\ON<3->{\B{\int \limits_{C_1} \FDR = \int \limits_{C_2} \FDR}}}$} \\
	&\\
	\hline
	\hline
\end{tabular}

\vspace{0.2cm}

\includegraphics[scale = 0.18]{Figures/10_2_Fig2.png}
\end{center}

\end{column}\begin{column}{0.0cm}\end{column}\end{columns}
\end{frame}

%~~~~~~~~~~~~~~~~~~~~~~~~~~~~~~~~~~~~~~~~~~~~~~~~~~~~~~~~~
\begin{frame}
\begin{columns}\begin{column}{12.0cm}

{III. Conservative Vector Fields and Potential Functions}

\vspace{0.1cm}

\scriptsize

\begin{center}

{\setlength{\extrarowheight}{-0.1cm}
\begin{tabular}{|p{0.4cm}p{10.6cm}|}
	\hline
	\hline
	&\\
	\DEF & A vector field ${\bf F}$ is {\bf conservative} if 
\ON<2->{\B{$\D{\WORK}$ is path independent.}}\\
	&\\
	\hline
	\hline
\end{tabular}}

\vspace{0.2cm}

{\setlength{\extrarowheight}{-0.1cm}
\begin{tabular}{|p{0.4cm}p{10.6cm}|}
	\hline
	\hline
	&\\
	\THM & If ${\bf F}$ is a conservative vector field then
	\ON<3->{\B{there exists a scalar function such that $\D{\bf F} = \nabla f$. }}\\
	&\\
	\hline
	\hline
\end{tabular}}

\vspace{0.2cm}

{\setlength{\extrarowheight}{-0.1cm}
\begin{tabular}{|p{0.4cm}p{10.6cm}|}
	\hline
	\hline
	&\\
	\DEF & The $f$ above is called a 
	\ON<4->{\B{{\bf potential for ${\bf F}$}.}}\\
	&\\
	\hline
	\hline
\end{tabular}}
\end{center}
{\setlength{\extrarowheight}{-0.2cm}
\begin{tabular}{|p{0.15cm}p{11.6cm}}
	\EX & Is $f(x,y,z)=x\sin (z^2) + e^{-y} + 13y$ a potential for $\D{{\bf F}(x,y,z) = \left[\sin(z^2),-e^{-y}+13,2xz\cos(z^2)\right]}$?\\
	& Is $\D{{\bf F}(x,y,z)}$ conservative?\\
	&\\
	& \ON<5->{\B{Since \quad $\nabla f = \left[\sin(z^2),-e^{-y}+13,2xz\cos(z^2)\right]$}} 
	\ON<6->{\B{$ = {\bf F}$}} \,\, 
	\ON<7->{\B{$ \implies \, f$ is a potential for ${\bf F}$}}\\
	&\\
	& \ON<8->{\B{${\bf F}$ is conservative since}}
	\ON<9->{\B{there exists an $f$ such that }}
	\ON<10->{\B{ $\nabla f = {\bf F}$}} \\
\end{tabular}}
 
\vspace{0.2cm}

{\setlength{\extrarowheight}{-0.2cm}
\begin{tabular}{|p{0.15cm}p{11.6cm}}
	\EX & Is $f(x,y,z)=xy\ln(y) + z$ a potential for $\D{{\bf F}(x,y,z) = \left[y\ln(y),\frac{x}{y},1\right]}$? \,\, Is $\D{{\bf F}(x,y,z)}$ conservative?\\
	&\\
	& \ON<11->{\B{Since \quad $\nabla f = [y\ln(y),x\ln(y)+x,1]$}} 
	\ON<12->{\B{$ \ne {\bf F}$}} \,\, 
	\ON<13->{\B{$ \implies \, f$ is NOT a potential for ${\bf F}$}}\\
	&\\
	& \ON<14->{\B{${\bf F}$ might or might not be conservative, we need more tools to determine this.}} \\
\end{tabular}}

\end{column}\begin{column}{0.0cm}\end{column}\end{columns}
\end{frame}

%~~~~~~~~~~~~~~~~~~~~~~~~~~~~~~~~~~~~~~~~~~~~~~~~~~~~~~~~~
\begin{frame}
\begin{columns}\begin{column}{12.0cm}

\scriptsize

\begin{center}
{\setlength{\extrarowheight}{-0.1cm}
\begin{tabular}{|p{0.4cm}p{10.6cm}|}
	\hline
	\hline
	&\\
	\THM & Let $C$ be a curve from $A$ to $B$.  If ${\bf F}$ is conservative (i.e. ${\bf F} = \nabla f$ for some $f$), then \\
	& \\
	\multicolumn{2}{|c|}{$\D{
	\ON<5->{\B{\int \limits_{C} \FDR }}
	\ON<6->{\B{\, = \int \limits_{C} \nabla f \cdot \, d{\bf r} }}
	\ON<7->{\B{\, = \left.f(x,y,z)\right|_A^B }}
	\ON<8->{\B{\, = f(B) - f(A)}}}$} \\
	&\\
	\hline
	\hline
\end{tabular}}
\end{center}
{\em Idea of Proof}:  
\ON<9->{\B{Let $\RT$ be a parameterization of $C$ with $t \in [a,b]$ then}}
%--------EQN----------
\begin{align*}
\ON<10->{\B{\int \limits_{C} \nabla f \cdot \, d{\bf r} = \int \limits_{a}^{b} \nabla f \cdot \RPT \, dt }} 
& \ON<11->{\B{= \int \limits_{a}^{b} \left[\PFPX,\PFPY,\PFPZ\right] \cdot \left[\DXDT,\DYDT,\DZDT\right] \, dt}} \\
& \ON<12->{\B{\, = \int \limits_{a}^{b} \left(\PFPX\DXDT + \PFPY\DYDT + \PFPZ\DZDT\right) \, dt}} \\
& \ON<13->{\B{\, = \int \limits_{a}^{b} \frac{df}{dt} \, dt}} \\
& \ON<14->{\B{\, = f(x(b),y(b),z(b)) - f(x(a),y(a),z(a))}} \\
& \ON<15->{\B{\, = f({\bf r}(b)) - f({\bf r}(a))}} \\
& \ON<16->{\B{\, = f(B) - f(A)}} \\
\end{align*}
%----------------------

\vspace{0.2cm}

\end{column}\begin{column}{0.0cm}\end{column}\end{columns}
\end{frame}

%~~~~~~~~~~~~~~~~~~~~~~~~~~~~~~~~~~~~~~~~~~~~~~~~~~~~~~~~~
\begin{frame}
\begin{columns}\begin{column}{12.0cm}

\vspace{0.1cm}

\scriptsize

There is a nice consequence of the previous theorem that is important in future discussions.

\vspace{0.4cm}

First a definition:

\vspace{0.4cm}

\DEF A curve, $C$, is {\bf closed} if it has 
\ON<2->{\B{\,\, no endpoints \,\, }} and 
\ON<3->{\B{\,\, encloses an area (not necessarily planar) \,\,\, }}.  Any point on the curve can be thought of as an
\ON<4->{\B{\,\, initial\,\, }} and an
\ON<5->{\B{\,\, end point of the curve.}}

\vspace{0.5cm}

\THM If ${\bf F}$ is {\bf conservative} (${\bf F} = \nabla f$) and $C$ is closed then

\vspace{0.1cm}

\begin{center}
	\ON<6->{\B{$\WORK = 0$}}
\end{center}

\vspace{0.2cm}

{\em Idea of Proof}:  
\ON<7->{\B{Take any point on $C$, lets call it $A$.}}  
\ON<8->{\B{Since $C$ is closed, it could be thought of as a curve from}}
\ON<9->{\B{ point $A$ to point $A$.}}

\vspace{0.3cm}

Then from the previous theorem, we have

\vspace{0.4cm}

\begin{center}
	$\D{
	\ON<10->{\B{\int \limits_{C} \FDR }}
	\ON<11->{\B{\, = \int \limits_{C} \nabla f \cdot \, d{\bf r} }}
	\ON<12->{\B{\, = \left.f(x,y,z)\right|_A^A }}
	\ON<13->{\B{\, = f(A) - f(A) = 0}}}$
\end{center}

\vspace{0.8cm}

\end{column}\begin{column}{0.0cm}\end{column}\end{columns}
\end{frame}

%~~~~~~~~~~~~~~~~~~~~~~~~~~~~~~~~~~~~~~~~~~~~~~~~~~~~~~~~~
\begin{frame}
\begin{columns}\begin{column}{12.0cm}

\vspace{0.1cm}

\scriptsize

\EX The vector field $\D{{\bf F} = [yz,xz,xy]}$ is conservative and has potential $f(x,y,z) = xyz$ (i.e. ${\bf F} = \nabla f)$.  Find the work done by {\bf F} on a mass along a curve, $C$, between $A(-1,3,9)$ and $B(1,6,-4)$. 

\vspace{0.2cm}

\SOL \ON<2->{\B{First note that since ${\bf F}$ is conservative, we have:}}

\vspace{0.2cm}

%--------EQN----------
\begin{align*}
	\ON<3->{\B{\WORK 
	    = \int \limits_{C} \nabla f \, d{\bf r} }}
	&\ON<4->{\B{\, = \int \limits_{C} \nabla[xyz] \, d{\bf r}}} \\
	&\ON<5->{\B{\, = xyz\biggr|_{(-1,3,9)}^{(1,6,-4)}}} \\
	&\ON<6->{\B{\, = (1)(6)(-4) - (-1)(3)(9)}} \\
	&\ON<7->{\B{\, = -24 + 27 = 3}}\\
\end{align*}
%----------------------

\vspace{3.2cm}

\end{column}\begin{column}{0.0cm}\end{column}\end{columns}
\end{frame}

%~~~~~~~~~~~~~~~~~~~~~~~~~~~~~~~~~~~~~~~~~~~~~~~~~~~~~~~~~
\begin{frame}
\begin{columns}\begin{column}{12.0cm}

{IV. Finding Potential Functions}

\vspace{0.3cm}

\scriptsize
If we are given a conservative vector field {\bf F}, then we know that there must be a corresponding potential function $f$ such that $\D{{\nabla f = {\bf F}}}$.

\vspace{0.4cm}

But what is $f$?

\vspace{0.5cm}

Let ${\bf F} = [M,N,P]$ \qquad (Note: $M, N, P$ are all functions of 
\ON<2->{\B{\,\, $x,y,$ and $z$.\,\,}})

\vspace{0.5cm}

Then
$\D{ 
\ON<3->{\B{\nabla f = {\bf F} \quad \implies \quad}}
\ON<4->{\B{\left[\PFPX,\PFPY,\PFPZ\right] = [M,N,P] \quad \implies \quad}}
\ON<5->{\B{\PFPX = M, \quad}}
\ON<6->{\B{\PFPY = N, \quad}}
\ON<7->{\B{\PFPZ = P}}
}$ 

\vspace{0.5cm}

So to find $f$, we must compute a sequence of 
\ON<8->{\B{partial anti-derivatives:}}

\vspace{0.5cm}

(i) 
\ON<9->{\B{\,\,\, \quad $\D{\PFPX = M} $ }}

\vspace{0.3cm}

(ii) 
\ON<10->{\B{\, \quad $\D{\PFPY = N}$ }}

\vspace{0.3cm}

(iii)
\ON<11->{\B{ \quad $\D{\PFPZ = P}$ }}

\vspace{1.2cm}

\end{column}\begin{column}{0.0cm}\end{column}\end{columns}
\end{frame}

%~~~~~~~~~~~~~~~~~~~~~~~~~~~~~~~~~~~~~~~~~~~~~~~~~~~~~~~~~
\begin{frame}
\begin{columns}\begin{column}{12.0cm}

\vspace{0.1cm}

\scriptsize
\EX Let $\D{{\bf F} = [e^x\cos y + yz,\, xz - e^x\sin y,\, xy + z]}$ \quad 

\vspace{0.2cm}

(a) Show ${\bf F} $ is conservative by finding an $f$ such that ${\bf F} = \nabla f$

\vspace{0.2cm}

$\D{
\ON<2->{\B{{\bf F} = \nabla f \quad \implies}} \quad 
\ON<3->{\B{
\left[\,
\only<4-15>{\UL}{\,e^x\cos y + yz\,},\, 
\only<16-35>{\UL}{\,xz - e^x\sin y\,},\, 
\only<36->{\UL}{\,xy + z\,}
\,\right] = 
\left[\,
\only<4-15>{\UL}{\,f_x\,},
\only<16-35>{\UL}{\,f_y\,},
\only<36->{\UL}{\,f_z\,}
\,\right]}}
}$

\vspace{0.2cm}
$\D{
	\setlength{\extrarowheight}{0.25cm}
	\begin{array}{rrcccll}
		\ON<8->{\B{\text{from above} \rightarrow}} 
		& & & 
		\ON<5->{\B{f_x}} & 
		\ON<6->{\B{=}} & 
		\ON<7->{\B{e^x\cos y + yz}} & 
		\\
		& & & 
		\ON<9->{\B{f}} & 
		\ON<10->{\B{=}} & 
		\ON<11->{\B{e^x\cos y + xyz + g(y,z)}} & 
		\ON<12->{\B{\leftarrow \,\, \text{anti-der wrt x}}}\\ 
		\ON<19->{\B{\text{from above} \rightarrow}} &
		\ON<18->{\B{\only<20->{\CAN}{xz} - \only<20->{\CAN}{e^x\sin y}}} & 
		\ON<17->{\B{=}} & 
		\ON<13->{\B{\, f_y}} & 
		\ON<14->{\B{=}} & 
		\ON<15->{\B{-\only<20->{\CAN}{e^x\sin y} + \only<20->{\CAN}{xz} + g_y}} & 
		\ON<16->{\B{\leftarrow \,\, \text{partial wrt y}}} \\
		& & & 
		\ON<21->{\B{g_y}} & 
		\ON<22->{\B{=}} & 
		\ON<23->{\B{0}} & 
		\ON<24->{\B{\leftarrow \,\, \text{solve for } g_y}} \\
		& & & 
		\ON<25->{\B{g}} & 
		\ON<26->{\B{=}} & 
		\ON<27->{\B{h(z)}} & 
		\ON<28->{\B{\leftarrow \,\, \text{anti-der wrt y}}}\\
		& & & 
		\ON<29->{\B{f}} & 
		\ON<30->{\B{=}} & 
		\ON<31->{\B{e^x\cos y + xyz + h(z)}} & 
		\ON<32->{\B{\leftarrow \,\, \text{plug $g$ into $f$}}} \\
		\ON<39->{\B{\text{from above} \rightarrow}} &
		\ON<38->{\B{\only<40->{\CAN}{xy} + z}} & 
		\ON<37->{\B{=}} & 
		\ON<33->{\B{\, f_z}} & 
		\ON<34->{\B{=}} & 
		\ON<35->{\B{\only<40->{\CAN}{xy} + h_z}} & 
		\ON<36->{\B{\leftarrow \,\, \text{partial wrt z}}} \\
		& & & 
		\ON<41->{\B{h_z}} & 
		\ON<42->{\B{=}} & 
		\ON<43->{\B{z}} & 
		\ON<44->{\B{\leftarrow \,\, \text{solve for } h_z}} \\
		& & &
		\ON<45->{\B{h}} & 
		\ON<46->{\B{=}} & 
		\ON<47->{\B{\frac{z^2}{2} + C}} & 
		\ON<48->{\B{\leftarrow \,\, \text{anti-der wrt z}}}\\
		& & & 
		\ON<49->{\B{f}} & 
		\ON<50->{\B{=}} & 
		\ON<51->{\B{e^x\cos y + xyz + \frac{z^2}{2} + C}} & 
		\ON<52->{\B{\leftarrow \,\, \text{plug $h$ into $f$}}} \\
		& & & 
		\ON<53->{\B{f}} & 
		\ON<53->{\B{=}} & 
		\ON<53->{\B{e^x\cos y + xyz + \frac{z^2}{2}}} & 
		\ON<54->{\B{\leftarrow \,\, \text{set $C=0$}}} \\
	\end{array}
}$

\vspace{7.5cm}

\end{column}\begin{column}{0.0cm}\end{column}\end{columns}
\end{frame}
%
%~~~~~~~~~~~~~~~~~~~~~~~~~~~~~~~~~~~~~~~~~~~~~~~~~~~~~~~~~
\begin{frame}
\begin{columns}\begin{column}{12.0cm}

\vspace{0.1cm}

\scriptsize

(b) Compute $\WORK$ \quad where $C$ is a curve from $A:(0,0,2)$ to $B:\left(2,\pi/2,0\right)$

\vspace{0.2cm}

\ON<2->{\B{Knowing \,\,  $\D{f}$ \, such that ${\bf F} = \GF$ \, makes this much easier since we don't need a parameterization for $C$. We don't even need to know what $C$ is.}}

\begin{center}
	$\D{
	\ON<3->{\B{\WORK }}
	\ON<4->{\B{\, = \int \limits_{A}^{B} \GF \, d{\bf r}}}
	\ON<5->{\B{ = f(B) - f(A)}}
	}$
\end{center}

\ON<6->{\B{and computing}}

\vspace{0.3cm}

\quad $\D{
\ON<7->{\B{f(B) }}
\ON<8->{\B{\, = e^{2}\cos(\pi/2) + (2)(\pi/2)(0) + \frac{(0)^2}{2}}}
\ON<9->{\B{ = 0}}}$

\quad $\D{
\ON<10->{\B{f(A) }}
\ON<11->{\B{\, = e^{0}\cos(0) + (0)(0)(2) + \frac{(2)^2}{2}}}
\ON<12->{\B{= 1+2=3}}}$

\vspace{0.3cm}

\ON<13->{\B{which gives}}

\vspace{0.3cm}

\quad $\D{
\ON<14->{\B{\WORK = 0 - 3 = -3}}}$

\vspace{7.5cm}
	
\end{column}\begin{column}{0.0cm}\end{column}\end{columns}
\end{frame}
%
%~~~~~~~~~~~~~~~~~~~~~~~~~~~~~~~~~~~~~~~~~~~~~~~~~~~~~~~~~
\begin{frame}
\begin{columns}\begin{column}{12.0cm}

\vspace{0.1cm}

\scriptsize
\EX Let $\D{{\bf F} = [y,\, x,\, 4]}$. \quad 

\vspace{0.2cm}

(a) Show ${\bf F}$ is conservative by finding an $f$ such that ${\bf F} = \nabla f$

\vspace{0.2cm}


$\D{
\ON<2->{\B{{\bf F} = \nabla f \quad \implies}} \quad 
\ON<3->{\B{
\left[\,
\only<4-15>{\UL}{\,y\,},\, 
\only<16-35>{\UL}{\,x\,},\, 
\only<36->{\UL}{\,4\,}
\,\right] = 
\left[\,
\only<4-15>{\UL}{\,f_x\,},
\only<16-35>{\UL}{\,f_y\,},
\only<36->{\UL}{\,f_z\,}
\,\right]}}
}$

\begin{center}
$\D{
	\setlength{\extrarowheight}{0.25cm}
	\begin{array}{rrcccll}
		\ON<8->{\B{\text{from above} \rightarrow}} 
		& & & 
		\ON<5->{\B{f_x}} & 
		\ON<6->{\B{=}} & 
		\ON<7->{\B{y}} & 
		\\
		& & & 
		\ON<9->{\B{f}} & 
		\ON<10->{\B{=}} & 
		\ON<11->{\B{xy + g(y,z)}} & 
		\ON<12->{\B{\leftarrow \,\, \text{anti-der wrt x}}}\\ 
		\ON<19->{\B{\text{from above} \rightarrow}} &
		\ON<18->{\B{\only<20->{\CAN}{x}}} & 
		\ON<17->{\B{=}} & 
		\ON<13->{\B{\, f_y}} & 
		\ON<14->{\B{=}} & 
		\ON<15->{\B{\only<20->{\CAN}{x} + g_y}} & 
		\ON<16->{\B{\leftarrow \,\, \text{partial wrt y}}} \\
		& & & 
		\ON<21->{\B{g_y}} & 
		\ON<22->{\B{=}} & 
		\ON<23->{\B{0}} & 
		\ON<24->{\B{\leftarrow \,\, \text{solve for } g_y}} \\
		& & & 
		\ON<25->{\B{g}} & 
		\ON<26->{\B{=}} & 
		\ON<27->{\B{h(z)}} & 
		\ON<28->{\B{\leftarrow \,\, \text{anti-der wrt y}}}\\
		& & & 
		\ON<29->{\B{f}} & 
		\ON<30->{\B{=}} & 
		\ON<31->{\B{xy + h(z)}} & 
		\ON<32->{\B{\leftarrow \,\, \text{plug $g$ into $f$}}} \\
		\ON<39->{\B{\text{from above} \rightarrow}} &
		\ON<38->{\B{4}} & 
		\ON<37->{\B{=}} & 
		\ON<33->{\B{\, f_z}} & 
		\ON<34->{\B{=}} & 
		\ON<35->{\B{0 + h_z}} & 
		\ON<36->{\B{\leftarrow \,\, \text{partial wrt z}}} \\
		& & & 
		\ON<41->{\B{h_z}} & 
		\ON<42->{\B{=}} & 
		\ON<43->{\B{4}} & 
		\ON<44->{\B{\leftarrow \,\, \text{solve for } h_z}} \\
		& & &
		\ON<45->{\B{h}} & 
		\ON<46->{\B{=}} & 
		\ON<47->{\B{4z + C}} & 
		\ON<48->{\B{\leftarrow \,\, \text{anti-der wrt z}}}\\
		& & & 
		\ON<49->{\B{f}} & 
		\ON<50->{\B{=}} & 
		\ON<51->{\B{xy + 4z + C}} & 
		\ON<52->{\B{\leftarrow \,\, \text{plug $h$ into $f$}}} \\
		& & & 
		\ON<53->{\B{f}} & 
		\ON<53->{\B{=}} & 
		\ON<53->{\B{xy + 4z}} & 
		\ON<54->{\B{\leftarrow \,\, \text{set $C=0$}}} \\
	\end{array}
}$
\end{center}

\end{column}\begin{column}{0.0cm}\end{column}\end{columns}
\end{frame}

%~~~~~~~~~~~~~~~~~~~~~~~~~~~~~~~~~~~~~~~~~~~~~~~~~~~~~~~~~
\begin{frame}
\begin{columns}\begin{column}{12.0cm}

\vspace{0.1cm}

\scriptsize

(b) Compute $\int \limits_{(1,1,1)}^{(2,3,-1)} y \, dx + x \, dy + 4 \, dz \quad 
\ON<2->{\B{\left( = \int \limits_{(1,1,1)}^{(2,3,-1)} \FDR  \quad where \quad {\bf F} = [y,x,4]\,\, \& \,\, d{\bf r} = [dx,dy,dz]\right)}}$

\vspace{0.2cm}

From part (a), we found that \, $\D{f(x,y,z) = 
\ON<3->{\B{xy + 4z}}}$

\vspace{0.2cm}

and therefore

%--------EQN----------
\begin{align*}
	%\ON<11->{\B{\int \limits_{(1,1,1)}^{(2,3,-1)} y \, dx + x \, dy + 4 \, dz}} 
	\ON<4->{\B{ \,         \int \limits_{(1,1,1)}^{(2,3,-1)} \FDR}}
	\ON<5->{\B{ \,\,  = \, \int \limits_{(1,1,1)}^{(2,3,-1)} \GF \cdot d{\bf r}}}
	&\ON<6->{\B{ \,\,  = \, f(2,3,-1) - f(1,1,1) }} \\
	&\ON<7->{\B{ \,\, = \, \left[(2)(3) + 4(-1)\right] - \left[(1)(1) + 4(1)\right] }} \\
	&\ON<8->{\B{\,\, = \, 2-5 = -3}}
\end{align*}
%----------------------

\vspace{7.5cm}
	
\end{column}\begin{column}{0.0cm}\end{column}\end{columns}
\end{frame}
%
%~~~~~~~~~~~~~~~~~~~~~~~~~~~~~~~~~~~~~~~~~~~~~~~~~~~~~~~~~
\begin{frame}
\begin{columns}\begin{column}{12.0cm}

{V. Conservative Test for Vector Fields}

\vspace{0.3cm}

\scriptsize

In general, to show a vector field is conservative by finding $f$ such that ${\bf F} = \nabla f$ can be time consuming and challenging.

\vspace{0.5cm}

Fortunately, there is a test to determine whether a vector field, {\bf F}, is conservative.

\vspace{0.5cm}

{\bf Conservative vector field test}: \quad  Let ${\bf F} = [M,N,P]$\, be a vector field.  

\vspace{0.5cm}

Assuming the partial derivatives of \, $M, N, P$\,  exist, then ${\bf F}$ is conservative if 

\begin{center}
	$\D{
	\ON<2->{\B{
	\frac{\PP P}{\PP y} = \frac{\PP N}{\PP z}, \quad 
	\frac{\PP M}{\PP z} = \frac{\PP P}{\PP x}, \quad 
	\frac{\PP N}{\PP x} = \frac{\PP M}{\PP y} \qquad 
	or }}
	\ON<3->{\B{\qquad P_y = N_z, \quad M_z=P_x, \quad N_x=M_y}}}$
\end{center}

This is also equivalent to showing

\begin{center}
	\ON<4->{\B{curl(${\bf F}$) = ${\bf 0}$}} \quad 
	\ON<5->{\B{or \quad $\nabla \times {\bf F} = {\bf 0}$}}
\end{center}

\EX Is $\D{{\bf F} = [2x-3,-z,\cos z]}$ conservative?

\vspace{0.2cm}
Since $\D{
\ON<6->{\B{\, P_y = (\cos z)_y}}
\ON<7->{\B{\,= 0}}
}$ \quad 
\ON<8->{\B{and}} \quad 
$\D{
\ON<9->{\B{N_z = (-z)_z}}
\ON<10->{\B{\, = -1}} \quad 
\ON<11->{\B{\implies}} \quad 
\ON<12->{\B{P_y \ne N_z}}}$

\vspace{0.2cm}

\ON<13->{\B{ So $\D{{\bf F}}$ is not conservative.}}

\vspace{1.2cm}

\end{column}\begin{column}{0.0cm}\end{column}\end{columns}
\end{frame}
%
%%%~~~~~~~~~~~~~~~~~~~~~~~~~~~~~~~~~~~~~~~~~~~~~~~~~~~~~~~~~
%%\begin{frame}
%%\frametitle{V. Differential Forms}
%%\scriptsize
%%\begin{columns}
%%\begin{column}{12.0cm}
%%\DEF $\FDR$ is known as a {\bf differential form} and 
%%
%%\begin{center}
%%	$\D{\FDR = \ON<3->{\B{[F_1,F_2,F_3] \cdot [dx,dy,dz] = \,}} \ON<4->{\B{F_1\,dx + F_2\,dy + F_2\,dz}}}$
%%\end{center}
%%
%%\vspace{1.2cm}
%%
%%\DEF If ${\bf F}$ is \ON<5->{\B{\,\, conservative\,\,}} then the differential form $\FDR$ is \ON<6->{\B{\,\,{\bf exact}.\,\,}}
%%
%%\vspace{1.2cm}
%%
%%Note: \quad \ON<7->{\B{$\FDR$ \, is exact}} \quad $\iff$ \quad \ON<8->{\B{${\bf F}$ \, is conservative}}	
%%\end{column}
%%\begin{column}{0.0cm}\end{column}
%%\end{columns}
%%\end{frame}
%
%~~~~~~~~~~~~~~~~~~~~~~~~~~~~~~~~~~~~~~~~~~~~~~~~~~~~~~~~~
\begin{frame}
\begin{columns}\begin{column}{12.0cm}

\vspace{0.1cm}

\scriptsize

\EX Let $\D{{\bf F} = [
\only<2->{\US{\B{ M }}}{\only<2->{\UB}{ 2xz }},\,
\only<2->{\US{\B{ N }}}{\only<2->{\UB}{ 3z }},\,
\only<2->{\US{\B{ P }}}{\only<2->{\UB}{ x^2+3y }}]}$. \quad Is ${\bf F}$ is conservative?

\vspace{0.2cm}

\ON<3->{\B{Check: }}

\vspace{0.5cm}

$\D{
\ON<4->{\B{P_y = (x^2+3y)_y }}
\ON<5->{\B{\,= 3
\qquad \text{and} \qquad }}
\ON<6->{\B{N_z = (3z)_z}}
\ON<7->{\B{\, = 3}}
\ON<8->{\B{\qquad \implies \qquad P_y=N_z \,\, \checkmark}}
}$

\vspace{0.5cm}

$\D{
\ON<9->{\B{M_z = (2xz)_z }}
\ON<10->{\B{\,= 2x
\qquad \text{and} \qquad }}
\ON<11->{\B{P_x = (x^2+3y)_x }}
\ON<12->{\B{\,= 2x }}
\ON<13->{\B{\qquad \implies \qquad M_z=P_x \,\, \checkmark}}
}$

\vspace{0.5cm}

$\D{
\ON<14->{\B{N_x = (3z)_x }}
\ON<15->{\B{\,= 0
\qquad \text{and} \qquad }}
\ON<16->{\B{M_y = (2xz)_y }}
\ON<17->{\B{\,= 0 }}
\ON<18->{\B{\qquad \implies \qquad N_x=M_y \,\, \checkmark}}
}$

\vspace{0.5cm}

\ON<19->{\B{Therefore, {\bf F} is conservative}}

\vspace{4.6cm}

\end{column}\begin{column}{0.0cm}\end{column}\end{columns}
\end{frame}

%~~~~~~~~~~~~~~~~~~~~~~~~~~~~~~~~~~~~~~~~~~~~~~~~~~~~~~~~~
\begin{frame}
\begin{columns}\begin{column}{12.0cm}

\vspace{0.1cm}

\scriptsize

\EX Let $\D{f = x^2z+3yz}$ \, and \, $\D{\bf F} = 2xz \VI + 3z \VJ + (x^2+3y) \VK$.

\vspace{0.2cm}

(a) \,\, Is \, $\D{f}$ \, the potential for \, $\D{\bf F}$? \quad 

\vspace{0.2cm}

\ON<2->{\B{Check:}}

\vspace{0.2cm}

$\D{
\ON<2->{\B{\GF = [(x^2z+3yz)_x, (x^2z+3yz)_y, (x^2z+3yz)_z] = [
\ON<3->{\B{\,\, 2xz \,\,}}, 
\ON<4->{\B{\,\, 3z \,\,}}, 
\ON<5->{\B{\,\, x^2+3y \,\,}}] = {\bf F} \,\, \checkmark}}
}$

\vspace{0.2cm}

\ON<2->{\B{Yes, $f$ is the potential for {\bf F}}}

\vspace{0.4cm}

(b) \,\, Suppose you know the work required to move a mass through ${\bf F}$ along a curve between the points $(1,1,1) \,\, \text{and} \,\, (d,2d,2)$ is 10 $N\cdot m$.  

\vspace{0.2cm}

What are the possible terminal points, $(d,2d,2)$, for $C$? 

\vspace{0.2cm}

\ON<6->{\B{Since {\bf F} is conservative, we have}}
%--------EQN----------
\begin{align*}
	\ON<8->{\B{10 = }}
	\ON<7->{\B{\WORK }}
	\ON<9->{\B{\, = \int \limits_{(d,2d,2)}^{(1,1,1)} \GF \cdot d{\bf r} }}&
	\ON<10->{\B{\, = f(d,2d,2) - f(1,1,1)}} \\ &
	\ON<11->{\B{\, = [(d)^2(2)+3(2d)(2)] - [(1)^2(1)+3(1)(1)]}} \\&
	\ON<12->{\B{\, = 2d^2 + 12d - 4}}
\end{align*}
%----------------------
$\D{
\ON<13->{\B{\quad 2d^2 + 12d - 4 = 10}}
}$ \\
$\D{
\ON<14->{\B{\quad 2d^2 + 12d - 14 = 0}}
}$ \\
$\D{
\ON<15->{\B{\quad 2(d^2 + 6d - 7) = 0}}
}$\\
$\D{
\ON<16->{\B{\quad 2(d-1)(d+7) = 0}} \quad 
\ON<17->{\B{\implies \quad d=1,-7}} \quad 
\ON<18->{\B{\implies \quad \text{Possible terminal points: 
\ON<19->{\B{$(1,2,2)$}} 
\ON<20->{\B{\,or $(-7,-14,2)$}}}}}
}$

\vspace{8.0cm}
	
\end{column}\begin{column}{0.0cm}\end{column}\end{columns}
\end{frame}
